% !TEX root = sga4.5-en.tex
% ENGLISH EDITION -- TRANSLATED FROM FRENCH (AI-assisted, needs human review).

\chapter{Derived categories}\label{VIII}
\blfootnote{by Jean-Louis Verdier}










\section{Triangulated categories: definitions and examples}\label{VIII:1}





\subsection{Definition of triangulated categories}\label{VIII:1-1}


\addtocounter{subsubsection}{-1}
\subsubsection{}\label{VIII:1-1-0}

Let $\cA$ be an additive category, $T$ an additive automorphism of $\cA$.
Given two objects $X$ and $Y$ of $\cA$, we set 
\[
  \hom^i(X,Y) = \hom_\cA(X,T^i(Y)) \qquad i\in \dZ 
\]
Let $X$, $Y$, $Z$ be three objects of $\cA$, $\alpha\in \hom^i(X,Y)$,
$\beta\in \hom^j(X,Y)$. We define the composite $\beta\circ\alpha$
in $\hom^{i+j}(X,Z)$ by: 
\[
  \beta\circ \alpha = \beta\circ T^i(\alpha) \text{.}
\]
One checks that this defines a new category whose groups of
morphisms between objects are graded. This new category will be
called, by abuse of language: the category $\cA$, graded by the
translation functor $T$.

Let $\cA$ be a category graded by a translation functor $T$.
When a morphism is named or denoted without specifying the degree, it
is always understood to be a morphism of degree zero.

Let $\cA$ and $\cA'$ be two additive categories graded by the functors
$T$ and $T'$, $F$ an additive functor from $\cA$ to $\cA'$. We say that $F$
\emph{is graded} if we are given an isomorphism of
functors\footnote{Text modified in 1976: the original text required an
equality $F\circ T=T'\circ F$. The case of functors in several variables
is discussed in \cite[XVIII 0.2]{sga4}. It raises a problem of signs.} 
\[
  F\circ T = T'\circ F 
\]

The definition of morphisms of graded functors is given in
\ref{VIII:4-1-1}.

$\cA$ being an additive category graded by the functor $T$, we
call triangle, a set of three objects $X$, $Y$, $Z$ and three
morphisms $u:X\to Y$, $v:Y\to Z$, $w:Z\to T(X)$ ($\deg w=1$). To denote
triangles we shall use the notation: $(X,Y,Z,u,v,w)$ or the diagram: 
\[\xymatrix{
  & Z \ar[dl]_-w \\
  X \ar[rr]^-u 
    & & Y \ar[lu]_-v 
}\qquad \deg(w)=1
\]
Let $(X,Y,Z,u,v,w)$, $(X',Y',Z',u',v',w')$ be two triangles of $\cA$; a
morphism from the first triangle to the second is a set of three
morphisms $f:X\to X'$, $g:Y\to Y'$, $h:Z\to Z'$, such that the following
diagram commutes: 
\[\xymatrix{
  X \ar[r]^-u \ar[d]^-f 
    & Y \ar[r]^-v \ar[d]^-g 
    & Z \ar[r]^-w \ar[d]^-h 
    & T(X) \ar[d]^-{T(f)} \\
  X' \ar[r]^-{u'} 
    & Y' \ar[r]^-{v'} 
    & Z' \ar[r]^-{w'} 
    & T(X') 
}\]





\subsubsection{}\label{VIII:1-1-1}

We call \emph{triangulated category} an additive category graded
by a translation functor, equipped with a family of triangles called
the family of distinguished triangles. This family must moreover satisfy
the axioms: 


\paragraph{TR1}
\phantomsection\hypertarget{VIII:TR1}{}

Every triangle isomorphic to a distinguished triangle is distinguished. Every
morphism $u:X\to Y$ is contained in a distinguished triangle
$(X,Y,Z,u,v,w)$. The triangle $(X,X,0,\operatorname{id}_{X'},0,0)$ is
distinguished. 


\paragraph{TR2}
\phantomsection\hypertarget{VIII:TR2}{}

For $(X,Y,Z,u,v,w)$ to be distinguished, it is necessary and sufficient that
$(Y,Z,T(X),v,w,-T(u))$ be distinguished. 


\paragraph{TR3}
\phantomsection\hypertarget{VIII:TR3}{}

Given distinguished $(X,Y,Z,u,v,w)$, $(X',Y',Z',u',v',w')$, for every
morphism $(f,g):u\to u'$, there exists a morphism $h:Z\to Z'$ such that $(f,g,h)$
is a morphism of triangles. 


\paragraph{TR4}
\phantomsection\hypertarget{VIII:TR4}{}

\[\xymatrix{
  & Z' \ar[dl]|{\deg=1} \\
  X \ar[rr]^-u 
    & & Y \ar[ul]_-i 
}\quad
\xymatrix{
  & X' \ar[dl]|{\deg(j)=1} \\
  Y \ar[rr]^-v 
    & & Z \ar[ul] 
}\quad
\xymatrix{
  & Y' \ar[dl]|{\deg=1} \\
  X \ar[rr]^-w 
    & & Z \ar[ul] 
}\]
being three distinguished triangles such that $w=v\circ u$, there exist two
morphisms $f:Z'\to Y'$, $g:Y' \to X'$ such that 
\begin{enumerate}
  \item $(\operatorname{id}_X, v, f)$ is a morphism of triangles
  \item $(u,\operatorname{id}_Z,g)$ is a morphism of triangles
  \item $(Z',Y',X',f,g,T(i) j)$ is a distinguished triangle
\end{enumerate}

Given two triangulated categories $\cC$ and $\cC'$, we call
\emph{exact functor} from $\cC$ to $\cC'$, any additive, graded functor
sending distinguished triangles to distinguished triangles. The
morphisms of exact functors are the morphisms of graded functors.

One deduces immediately from axioms \hyperlink{VIII:TR1}{TR1},
\hyperlink{VIII:TR2}{TR2}, \hyperlink{VIII:TR3}{TR3} the following property: 





\begin{proposition}\label{VIII:1-1-2}
Let $(X,Y,Z,u,v,w)$ be a distinguished triangle and $M$ an object of a
triangulated category. We then have the exact sequences: 
\[\xymatrix{
  \cdots \ar[r] 
    & \hom^i(M,X) \ar[r]^-{\hom(M,u)} 
    & \hom^i(M,Y) \ar[r]^-{\hom(M,v)} 
    & \hom^i(M,Z) \ar[r]^-{\hom(M,w)} 
    & \hom^{i+1}(M,X) \ar[r] 
    & \cdots
}\]
\[\xymatrix{
  \cdots \ar[r] 
    & \hom^i(Z,M) \ar[r]^-{\hom(v,M)} 
    & \hom^i(Y,M) \ar[r]^-{\hom(u,M)} 
    & \hom^i(X,M) \ar[r]^-{\hom(w,M)} 
    & \hom^{i+1}(Z,M) \ar[r] 
    & \cdots
}\]
\end{proposition}

One then deduces from this proposition statements of the form:
\begin{itemize}
  \item $(f,g,h)$ being a morphism of distinguished triangles and $f,g$
    isomorphisms.
  \item The distinguished triangles built on a morphism
    (\hyperlink{VIII:TR1}{TR1}), are all isomorphic (\emph{but the
    isomorphisms are not, in general, uniquely determined}).
  \item When in a distinguished triangle, one of the morphisms admits a kernel
    or a cokernel, it splits i.e. it is of the form: 
    \[\xymatrix{
      & Z+T(X) \ar[dl] \\
      X+Y \ar[rr] 
        & & Y+Z \ar[ul]
    }\]
  \item The sum of two distinguished triangles is distinguished. (One uses
    axiom \hyperlink{VIII:TR4}{TR4}). 
\end{itemize}










\subsection{Examples}\label{VIII:1-2}





\subsubsection{}\label{VIII:1-2-1}

We first fix some notation. Let $\cA$ be an additive
category $\eC(\cA)$ will denote the following category:
\begin{itemize}
  \item The objects of $\eC(\cA)$ will be the complexes of $\cA$, with no
    bound on the degree, with differential of degree $+1$.
  \item The morphisms of $\eC(\cA)$ will be the morphisms of complexes (which
    commute with the differential) preserving the degree.
\end{itemize}

Let $T$ be the following functor: for every object $X^\bullet$ of $\eC(\cA)$ 
\begin{align*}
  T(X^\bullet)_i &= (X^\bullet)_{i+1} \\
  d_i(T(X^\bullet)) &= -d_{i+1}(X^\bullet)  \text{.}
\end{align*}
On morphisms the functor $T$ acts as follows: for every
morphism $f:X^\bullet \to Y^\bullet$, the morphism
$T(f):T(X^\bullet) \to T(Y^\bullet)$ is defined by 
\[
  T(f)_i = f_{i+1} \text{.}
\]
One checks immediately that $T$ is an additive automorphism of
$\eC(\cA)$. The morphisms of degree $n$, defined by means of this
translation functor, are then the morphisms which raise the degree by $n$, and
which commute or anticommute with the differential according to the parity of
$n$. One thus recovers the generally adopted definition.

$\sC$ will denote the following family of triangles: a triangle
$(X^\bullet,Y^\bullet,Z^\bullet,u,v,w)$ is an element of the family
if\footnote{Text modified in 1976, to ensure the validity of
\ref{VIII:1-2-4}.}
\begin{itemize}
  \item $Z^\bullet$ is the simple complex associated to the double complex 
    \[\xymatrix{
      \cdots \ar[r] 
        & 0 \ar[r] 
        & X^\bullet \ar[r]^-u 
        & Y^\bullet \ar[r] 
        & 0 \ar[r] 
        & \cdots 
    }\]
    where the objects of $Y^\bullet$ are the objects of first degree zero,
    and the objects of $X^\bullet$, the objects of first degree $-1$.
  \item $v$ is the image under the functor: double complex $\mapsto$ simple
    complex, of the morphism of double complexes: 
    \[\xymatrix{
      \cdots \ar[r] 
        & 0 \ar[r] \ar[d] 
        & 0^\bullet \ar[r] \ar[d] 
        & Y^\bullet \ar[r] \ar[d] 
        & 0 \ar[r] \ar[d] 
        & \cdots \\
      \cdots \ar[r] 
        & 0 \ar[r] 
        & X^\bullet \ar[r] 
        & Y^\bullet \ar[r] 
        & 0 \ar[r] 
        & \cdots 
    }\]
  \item $w$ is the opposite of the image, under the same functor, of the morphism of
    double complexes: 
    \[\xymatrix{
      \cdots \ar[r] 
        & 0 \ar[r] 
        & X^\bullet \ar[r] 
        & 0 \ar[r] 
        & 0 \ar[r] 
        & \cdots \\
      \cdots  \ar[r] 
        & 0 \ar[r] \ar[u] 
        & X^\bullet \ar[r] \ar[u] 
        & Y^\bullet \ar[r] \ar[u] 
        & 0 \ar[r] \ar[u] 
        & \cdots 
    }\]
\end{itemize}





\subsubsection{}\label{VIII:1-2-2}

The set of morphisms homotopic to zero is a two-sided ideal (if
$f$ and $g$ are composable and if $f$ or $g$ is homotopic to zero, the composite
is homotopic to zero. If $f$ and $g$ are homotopic to zero, the morphism
$f\oplus g$ is homotopic to zero). For every pair of objects $X^\bullet$,
$Y^\bullet$ we denote by $\operatorname{Htp}(X^\bullet,Y^\bullet)$ the
subgroup of $\hom_{\eC(\cA)}(X^\bullet,Y^\bullet)$ of morphisms homotopic
to zero. $\eK(\cA)$ will then be the category whose objects are the
objects of $\eC(\cA)$ and whose morphisms with source $X^\bullet$ and target
$Y^\bullet$ form the group: 
\[
  \hom_{\eC(\cA)}(X^\bullet,Y^\bullet) / \operatorname{Htp}(X^\bullet,Y^\bullet) \text{.}
\]
The composition law on morphisms of $\eK(\cA)$ is defined from
the composition law on morphisms of $\eC(\cA)$ by passing to the
quotient. $\eK(\cA)$ is an additive category. The translation functor
$T$ on $\eC(\cA)$ passes to the quotient (the translate of a morphism homotopic
to zero, is homotopic to zero). The resulting functor is an additive
automorphism of $\eK(\cA)$ which we still denote by $T$. By abuse of notation
$\eK(\cA)$ will henceforth denote the category $\eK(\cA)$ graded by the
functor $T$.

We call distinguished triangle in $\eK(\cA)$, any triangle
\emph{isomorphic} to a triangle coming from $\sC$. The family of distinguished
triangles satisfies axioms \hyperlink{VIII:TR1}{TR1},
\hyperlink{VIII:TR2}{TR2}, \hyperlink{VIII:TR3}{TR3} and
\hyperlink{VIII:TR4}{TR4}. By abuse of notation $\eK(\cA)$ will denote the
category $\eK(\cA)$ triangulated by the family of triangles defined
above. 





\subsubsection{}\label{VIII:1-2-3}

We further define a whole arsenal of subcategories:
\begin{itemize}
  \item A complex is said to be bounded below if all its objects of
    negative degree, except at most finitely many, are zero. One defines
    likewise the complexes bounded above, the bounded complexes. We
    denote by $\eK^+(\cA)$, $\eK^-(\cA)$, $\eK^b(\cA)$ the
    full subcategories of $\eK(\cA)$ generated by these families
    of objects. These subcategories are ``\emph{triangulated
    subcategories}'' of $\eK(\cA)$: every distinguished triangle two of whose
    objects are objects of the subcategory, is isomorphic to a
    triangle all three of whose objects are objects of the subcategory.
  \item Suppose that $\cA$ is an abelian category. A complex is
    said to be cohomologically bounded below if all its cohomology objects
    of negative degree vanish except finitely many of
    them. One defines analogously the cohomologically
    bounded above complexes, the cohomologically bounded complexes, the
    acyclic complexes. $\eK^{\infty,+}(\cA)$, $\eK^{\infty,-}(\cA)$,
    $\eK^{\infty,b}(\cA)$, $\eK^{\infty,\varnothing}(\cA)$ will denote the
    full subcategories generated by these families of objects. They are
    triangulated subcategories.

    One also gets triangulated subcategories $\eK^{+,b}(\cA)$,
    $\eK^{+,\varnothing}(\cA)$, \ldots etc \ldots (the first sign in
    the exponent gives information on the objects of the complexes, the
    second sign on the cohomology objects).
  \item $\cA$ no longer necessarily abelian, let $O$ be a
    set of objects of $\cA$ stable under isomorphism and direct sum.
    $O(\cA)$ will denote the full subcategory of $\eC(\cA)$ generated
    by the complexes whose objects in every degree are objects of $O$.
    We denote by $\underline O(\cA)$ the corresponding triangulated
    subcategory in $\eK(\cA)$. When $\cA$ is abelian, one may
    take for the set $O$, the set $I$ of injectives or the set $P$
    of projectives.
\end{itemize}





\subsubsection{}\label{VIII:1-2-4}

$\eK(\cA)$ is uniquely determined by the data of $\eC(\cA)$ and the
family $\sC$. More precisely: every graded additive functor from
$\eC(\cA)$ to a triangulated category, sending every triangle of
$\sC$ to a distinguished triangle, factors uniquely through
$\eK(\cA)$, and the resulting functor is exact.

For every complex $X^\bullet$ of $\eC(\cA)$, $\bar X^\bullet$ will denote the
complex obtained from $X^\bullet$ by killing the differential. A
three-term sequence: 
\[\xymatrix{
  0 \ar[r] 
    & X^\bullet \ar[r] 
    & Y^\bullet \ar[r] 
    & Z^\bullet \ar[r] 
    & 0 
}\]
will be called a semi-split sequence if the sequence 
\[\xymatrix{
  0 \ar[r] 
    & \bar X^\bullet \ar[r] 
    & \bar Y^\bullet \ar[r] 
    & \bar Z^\bullet \ar[r] 
    & 0 
}\]
is split; i.e. if it is isomorphic to the sequence: 
\[\xymatrix{
  0 \ar[r] 
    & \bar X^\bullet \ar[r] 
    & \bar X^\bullet + \bar Z^\bullet \ar[r] 
    & \bar Z^\bullet \ar[r] 
    & 0 \text{.}
}\]

Let $0 \to X^\bullet \to Y^\bullet \to Z^\bullet \to 0$ be a
semi-split sequence. Choose a splitting i.e. for every $n$ an isomorphism
\[
  (Y^\bullet)^n \simeq (X^\bullet)^n + (Z^\bullet)^n \text{.}
\]
The matrix of the differential of $Y^\bullet$, in the above direct sum
decomposition is of the form: 
\[
  \begin{pmatrix}
    d_{X^\bullet} & \varphi_n \\
    0 & d_{Y^\bullet} 
  \end{pmatrix}
\]
where $\varphi_n$ is a morphism from $(Z^\bullet)^n$ to
$(X^\bullet)^{n+1}$. The $\varphi_n$ determine a morphism of complexes
\[
  \varphi:Z^\bullet \to T(X^\bullet) \text{.}
\]
When one changes the splitting, the class of $\varphi$ modulo homotopy does
not change. Moreover, denoting by $\underline u$, $\underline v$,
$\underline\varphi$ the images in $\eK(\cA)$ of the morphisms $u$, $v$,
$\varphi$, the triangle
$(X^\bullet,Y^\bullet,Z^\bullet,\underline u,\underline v,\underline\varphi)$
is distinguished. Denote by $\mathsf{SSS}(\cA)$ the category of
semi-split sequences of $\eC(\cA)$, and by $\mathsf{TrK}(\cA)$ the
category of distinguished triangles of $\eK(\cA)$. What precedes
allows us to define a functor:
\[
  \rho:\mathsf{SSS}(\cA) \to \mathsf{TrK}(\cA) \text{.}
\]
This functor is essentially surjective. 










\subsection{Examples of cohomological functors and exact functors}\label{VIII:1-3}





\begin{definition}
Let $\cA$ be a triangulated category, $\cB$ an abelian category.
An additive functor $F:\cA\to \cB$ is called a \emph{cohomological functor} if
for every distinguished triangle $(X,Y,Z,u,v,w)$ the sequence 
\[\xymatrix{
  F(X) \ar[r]^-{F(u)} 
    & F(Y) \ar[r]^-{F(v)} 
    & F(Z) 
}\]
is exact.
\end{definition}

The functor $F\circ T^i$ will often be denoted $F^i$. By virtue of axiom
\hyperlink{VIII:TR2}{TR2} for triangulated categories, we have the
unbounded exact sequence: 
\[\xymatrix{
  \cdots \ar[r] 
    & F^i(X) \ar[r] 
    & F^i(Y) \ar[r] 
    & F^i(Z) \ar[r] 
    & F^{i+-}(X) \ar[r] 
    & \cdots 
}\]





\subsubsection{Examples}\label{VIII:1-3-2}

\begin{enumerate}
  \item Let $\cA$ be a triangulated category. For every object $X$ of $\cA$
    the functor $\hom_\cA(X,-)$ is a cohomological functor. The functor
    $\hom_\cA(-,X)$ is a cohomological functor $\cA^\circ\to\mathsf{Ab}$.
  \item Let $\cA$ be an abelian category, $\eC(\cA)$ the category of
    complexes of $\cA$. We denote by $\h^0:\eC(\cA)\to\cA$ the following
    functor: let $X^\bullet$ be an object of $\eC(\cA)$. We set
    \[
      \h^0(X^\bullet) = \ker d^0 / \im{d^{-1}} \text{.}
    \]
    $\h^0$ kills morphisms homotopic to zero, hence factors uniquely
    through $\eK(\cA)$. We still denote by
    $\h^0:\eK(\cA) \to \cA$ the resulting functor.
\end{enumerate}





\subsubsection{Example of an exact bifunctor}\label{VIII:1-3-3}

Let $\cA$, $\cA'$, $\cA''$ be three additive categories,
\[
  P:\cA\times \cA'\to \cA''
\]
a bilinear functor (i.e. additive in each argument
separately). One deduces the bilinear functor:
\[
  P^\ast:\eC(\cA)\times \eC(\cA') \to \eC(\cA'')
\]
as follows:

Let $X^\bullet$ be an object of $\eC(\cA)$ and $Y^\bullet$ an object of
$\eC(\cA')$. $P(X^\bullet,Y^\bullet)$ is a double complex of $\cA''$. We
then set: $P^\ast(X^\bullet,Y^\bullet) = $ simple complex associated to
$P(X^\bullet,Y^\bullet)$.

Let $f$ be a morphism of $\eC(\cA)$ (resp. $\eC(\cA')$) homotopic to zero
and $Z^\bullet$ an object of $\eC(\cA')$ (resp. $\eC(\cA)$). The morphism
$P^\ast(f,Z^\bullet)$ (resp. $P^\bullet(Z^\bullet,f)$) is then homotopic to
zero. One deduces that $P^\ast$ defines in a unique way a
functor:
\[
  P^\bullet:\eK(\cA) \times \eK(\cA') \to \eK(\cA'') \text{.}
\]
$P^\bullet$ is an exact bifunctor.

In particular, let $\cA$ be an additive category. One may take for $P$ the
functor: 
\begin{align*}
  \cA^\circ\times \cA &\to \mathsf{Ab} \\
  (X,Y) &\mapsto \hom(X,Y) 
\end{align*}

One then obtains by the preceding construction a functor
\[
  \hom^\bullet:\eK(\cA)^\circ \times \eK(\cA) \to \eK(\mathsf{Ab})
\]
which, composed with the functor $\h^0:\eK(\mathsf{Ab})\to \mathsf{Ab}$,
clearly gives back the functor $\hom_{\eK(\cA)}$. 




















\section{Quotient categories}\label{VIII:2}










\subsection{Thick subcategories. Multiplicative systems of morphisms}\label{VIII:2-1}

We begin by giving two definitions. 





\begin{definition}\label{VIII:2-1-1}
A subcategory $\cB$ of a triangulated category $\cA$ is said to be
\emph{thick} if $\cB$ is a full triangulated subcategory of $\cA$
and moreover $\cB$ has the following property:

For every morphism $f:X\to Y$, factoring through an object of $\cB$ and contained
in a distinguished triangle $(X,Y,Z,f,g,h)$, where $Z$ is an object of $\cB$,
the source of $f$ and the target of $f$ are objects of $\cB$.
\end{definition}

The set of thick subcategories of $\cA$ will be denoted $\sN$.

The intersection of any family of thick subcategories is
a thick subcategory. 





\begin{definition}
Let $\cA$ be a triangulated category. A set $S$ of morphisms of $\cA$
is called a \emph{multiplicative system of morphisms} if it has the
five following properties: 

\paragraph{FR1}
\phantomsection\hypertarget{VIII:FR1}{}
If $f,g\in S$ and if $f$ and $g$ are composable, $f\circ g\in S$. For every
object $X$ of $\cA$, the identity morphism of $X$ is an element of $S$. 

\paragraph{FR2}
\phantomsection\hypertarget{VIII:FR2}{}
In the category $\cA$, every diagram 
\[\xymatrix{
  & Y \ar[d]^-{s\in S} \\
  Z \ar[r]^-f 
    & X 
}\]
completes to a commutative diagram 
\[\xymatrix{
  P \ar[r]^-g \ar[d]^-{t\in S} 
    & Y \ar[d]^-{s\in S} \\
  Z \ar[r]^-f 
    & X 
}\]
Moreover the symmetric property holds. 

\paragraph{FR3}
\phantomsection\hypertarget{VIII:FR3}{}
If $f$ and $g$ are morphisms, the following properties are
equivalent:
\begin{enumerate}[\indent i)]
  \item there exists $s\in S$ such that $s\circ f=s\circ g$
  \item there exists $t\in S$ such that $f\circ t=g\circ t$
\end{enumerate}

\paragraph{FR4}
\phantomsection\hypertarget{VIII:FR4}{}
Let $T$ be the translation functor of $\cA$. For every element $s$ of $S$,
$T(s)\in S$. 

\paragraph{FR5}
\phantomsection\hypertarget{VIII:FR5}{}
Given two distinguished triangles $(X,Y,Z,u,v,w)$, $(X',Y',Z',u',v',w')$,
$f$ and $g$ two elements of $S$ such that the pair $(f,g)$ is a
morphism from $u$ to $u'$, there exists a morphism $h\in S$, such that
$(f,g,h)$ is a morphism of triangles. 
\end{definition}

A multiplicative system of morphisms is said to be \emph{saturated} if it has
the following property:

A morphism $f$ belongs to $S$ if and only if there exist two
morphisms $g$ and $g'$ such that $g\circ f\in S$ and $f\circ g'\in S$.

The intersection of any family of saturated multiplicative systems is
a saturated multiplicative system.

The set of saturated multiplicative systems will be denoted $\sS$.

Let $\cB$ be a thick subcategory; we denote by $\varphi(\cB)$
the set of morphisms $f$ contained in a distinguished triangle
$(X,Y,Z,f,g,h)$ where $Z$ is an object of $\cB$. $\varphi(\cB)$ is a
saturated multiplicative system.

Let $S$ be a saturated multiplicative system; we denote by $\psi(S)$ the
full subcategory generated by the objects $Z$ contained in a
distinguished triangle $(X,Y,Z,f,g,h)$ where $f$ is an element of $S$. The
subcategory $\psi(S)$ is a thick subcategory.

The main result of this section is then the following: $\varphi$ is an
isomorphism (preserving the order relation defined by inclusion) from
$\sN$ onto $\sS$. The inverse isomorphism is none other than $\psi$. 










\subsection{Universal problems}\label{VIII:2-2}

Let $\cA$ and $\cA'$ be two triangulated categories and $F$ an exact
functor from $\cA$ to $\cA'$. Let $S(F)$ be the set of morphisms of $\cA$
sent by $F$ to isomorphisms and $\cB(F)$ the full subcategory
generated by the objects of $\cA$ sent by $F$ to
zero objects of $\cA'$. $S(F)$ is a saturated multiplicative system and
$\cB(F)$ is a thick subcategory. Moreover $S(F) = \varphi(\cB(F))$.
Let then $\cA$ be a triangulated category, $\cB$ a thick subcategory
and $S=\varphi(\cB)$ the corresponding saturated multiplicative system.
The two universal problems below are equivalent.


\paragraph{Problem 1}
\phantomsection\hypertarget{VIII:prob1}{}
Find a triangulated category $\cA/\cB$ and an exact functor
$Q:\cA\to \cA/\cB$ such that every exact functor from $\cA$ to a
triangulated category $\cA'$, sending the objects of $\cB$ to zero objects of
$\cA'$, admits, uniquely, a factorisation of the form $G\circ Q$
where $G$ is an exact functor. 


\paragraph{Problem 2}
\phantomsection\hypertarget{VIII:prob2}{}
Find a triangulated category $\cA_S$ and an exact functor
$Q:\cA\to \cA_S$ such that every exact functor from $\cA$ to a
triangulated category $\cA'$, sending the morphisms of $S$ to isomorphisms of
$\cA'$, admits, uniquely, a factorisation of the form
$G\circ Q$ where $G$ is an exact functor.

(The author begs the reader to excuse him for the use of this
old-fashioned and retrograde language.) We shall show, in the next
section, that \hyperlink{VIII:prob2}{problem 2} admits a solution. Hence
the corresponding \hyperlink{VIII:prob1}{problem 1} admits the same solution. 










\subsection{Calculus of fractions}\label{VIII:2-3}





\subsubsection{}\label{VIII:2-3-1}

Let $\cA$ be a triangulated category, $S$ a saturated
multiplicative system. Denote by $\cA[S^{-1}]$ the fraction category of $\cA$
for the set of morphisms $S$ and $Q$ the canonical functor from $\cA$ to
$\cA[S^{-1}]$. The pair $(\cA[S^{-1}],Q)$ solves the following problem: every
functor from $\cA$ to any category (not necessarily
additive) sending morphisms of $S$ to isomorphisms, factors, uniquely,
through $(\cA[S^{-1}],Q)$. Such a pair always exists, with no
hypothesis on the set of morphisms $S$ \cite[1.1]{gz67}\footnote{The
original text references [C.G.G] here. Most likely, this is the lost work
\emph{Catégories et foncteurs} by Chevalley, Gabriel and Grothendieck,
referenced in\cite{sc72}. }. However the set
of morphisms $S$, being a saturated multiplicative system, has the
properties \hyperlink{VIII:FR1}{FR1}, \hyperlink{VIII:FR2}{FR2},
\hyperlink{VIII:FR3}{FR3}. Consequently, by \cite{gz67}, the category
$\cA[S^{-1}]$ can be obtained by a calculus of fractions, on the right or on the
left. We recall how $\cA[S^{-1}]$ is obtained by a right
calculus of fractions. (The left calculus of fractions is obtained by
reversing the arrows.) 





\subsubsection{}\label{VIII:2-3-2}

The objects of $\cA[S^{-1}]$ are the objects of $\cA$.
\begin{itemize}
  \item For every object $X$ of $\cA$, denote by $S_X$ the category of
    arrows in $S$, with target $X$. To every object $Y$ of
    $\cA$, we associate, functorially in $Y$, the following functor on
    $S_X^\circ$ with values in the category of sets:
    ($S_X^\circ$ denotes the opposite category) 
    \[
      H_Y:s\mapsto \hom(\operatorname{source}(s),Y) \text{.}
    \]
    We then set:
    \[
      \hom_{\cA[S^{-1}]}(X,Y) = \varinjlim_{S_X^\circ} H_Y \text{.}
    \]
    For everything concerning inductive limits, we refer to
    \cite{ar62}.

    One then observes, with pleasure, using \hyperlink{VIII:FR1}{FR1},
    \hyperlink{VIII:FR2}{FR2}, \hyperlink{VIII:FR3}{FR3} that the category
    $S_X^\circ$ has properties L1, L2, L3. Inductive limits
    therefore have all the good properties of inductive limits
    over filtered ordered sets.
  \item Let $X$, $Y$, $Z$, be three objects of $\cA$, $a$ an element of
    $\hom_{\cA[S^{-1}]}(X,Y)$ and $b$ an element of
    $\hom_{\cA[S^{-1}]}(Y,Z)$. Let $s$ be an object of $S_X$, $\underline a$ an
    element of $\hom(\operatorname{source}(s),Y)$ whose image is $a$ and
    an object of $S_y$, $\underline b$ an element of
    $\hom(\operatorname{source}(t),Z)$ whose image is $b$. We then have a
    diagram: 
    \begin{equation*}\tag{1}\label{VIII:eq:2-3-2-1}
    \xymatrix{
      & \bullet \ar[dl]_-s \ar[dr]^-{\underline a} 
        & 
        & \bullet \ar[dl]_-t \ar[dr]^-{\underline b} \\
      X 
        & 
        & Y 
        & 
        & Z 
    }
    \end{equation*}
    which, by \hyperlink{VIII:FR2}{FR2}, can be completed to a
    diagram: 
    \[\xymatrix{
      & & \bullet \ar[dl]_-{t'} \ar[dr]^-{\underline c} \\
      & \bullet \ar[dl]_-s \ar[dr]^-{\underline a} 
        & & \bullet \ar[dl]_-t \ar[dr]^-{\underline b} \\
      X 
        & & Y 
        & & Z 
    }\]
    Denote by $b\circ a$ the image in $\hom_{\cA[S^{-1}]}(X,Z)$ of the morphism
    $\underline{b\circ c}$. One checks that thanks to properties
    \hyperlink{VIII:FR1}{FR1}, \hyperlink{VIII:FR2}{FR2},
    \hyperlink{VIII:FR3}{FR3}, $b\circ a$ does not depend on the chosen representatives
    $\underline a$ and $\underline b$ and that it does not depend either
    on the way of completing diagram \eqref{VIII:eq:2-3-2-1}. One
    further checks that this defines a category $\cA[S^{-1}]$.
\end{itemize}

We denote by $Q$ the obvious functor: $\cA\to \cA[S^{-1}]$.

$\cA$ being an additive category, $\cA[S^{-1}]$ is an additive
category. Moreover, one shows using \hyperlink{VIII:FR4}{FR4} that there
exists one and only one translation functor on $\cA[S^{-1}]$, which we
denote $T$, satisfying the relation:
\[
  Q\circ T = T\circ Q \text{.}
\]
Finally, using \hyperlink{VIII:FR5}{FR5}, one shows that there exists one
and only one triangulated structure on $\cA[S^{-1}]$ such that the functor $Q$
is exact. The distinguished triangles of $\cA[S^{-1}]$ are none other than the
triangles isomorphic to the images, under $Q$, of the distinguished triangles of $\cA$.
Denoting by $\cA_S$, the category $\cA[S^{-1}]$ equipped with this
triangulated structure, one shows without difficulty that the pair
$(\cA_S,Q)$ is a solution to \hyperlink{VIII:prob2}{problem 2}. 





\begin{definition}\label{VIII:2-3-3}
Let $\cA$ be a triangulated category, $\cB$ a thick subcategory
of $\cA$, $(Q,\cA/\cB)$ the solution to \hyperlink{VIII:prob1}{problem 1}
(\S \ref{VIII:2-2}). $\cA/\cB$ will be called the \emph{quotient category} of
$\cA$ by $\cB$. $Q$ will be called the canonical quotient
functor. More generally let $\cN$ be a triangulated subcategory of
$\cA$ and let $\underline\cN$ be the smallest thick subcategory
containing $\cN$. The quotient category of $\cA$ by $\cN$: $\cA/\cN$ will by
definition be $\cA/\underline\cN$. 
\end{definition}










\subsection{Properties of quotient categories}\label{VIII:2-4}





\subsubsection{}\label{VIII:2-4-1}

Let $\cA$ be a triangulated category and $H$ a cohomological functor with
values in an abelian category $\cG$. Let $\cB_H$ be the
full subcategory generated by the objects all of whose translates
are sent by $H$ to zero objects of $\cG$ and $S_H$ the set of
morphisms all of whose translates are sent by $H$ to
isomorphisms of $\cG$. $\cB_H$ is a thick subcategory of $\cA$;
$S_H$ is a saturated multiplicative system. Moreover $S_H=\varphi(\cB_H)$
(\S \ref{VIII:2-1}). The functor $H$ factors uniquely through
$(Q,\cA/\cB_H)$. 





\begin{theorem}\label{VIII:2-4-2}
Let $\cA$ be a triangulated category, $\cB$ a triangulated
subcategory, $\cN$ a thick subcategory of $\cA$, $S=\varphi(\cN)$
the corresponding saturated multiplicative system.
\begin{enumerate}[(a)]
  \item The category $\cN\cap\cB$ is a thick subcategory of
    $\cB$. The corresponding multiplicative system is $S\cap \cB$.
  \item The two following properties are equivalent:
    \begin{enumerate}[(i)]
      \item For every object $X$ of $\cB$ and every morphism $s:R\to X$ where $s$
        is an element of $S$, there exists a morphism $s':R'\to R$ where
        $s\circ'$ is an element of $S\cap \cB$; and
        $R'\in\operatorname{Ob}\cB$.
      \item Every morphism from an object $X$ of $\cB$ to an object $Y$ of $\cN$
        factors through an object of $\cN\cap \cB$.
    \end{enumerate}
  \item Properties (i)' and (ii)' obtained from (i) and (ii) by
    reversing the arrows are equivalent.
  \item If properties (i) and (ii) hold or if
    properties (i)' and (ii)' hold, the canonical functor
    \[\xymatrix{
      \cB/\cN\cap \cB \ar[r] & \cA/\cN
    }\]
    is faithful.

    Since this functor is injective on objects, it realises
    $\cB/\cN\cap\cB$ as a subcategory of $\cA/\cN$.
  \item If moreover $\cB$ is a full subcategory of $\cA$,
    $\cB/\cN\cap \cB$ is a full subcategory of $\cA/\cN$. 
\end{enumerate}
\end{theorem}





\begin{corollary}\label{VIII:2-4-3}
Let $\cA\subset \cB\subset \cC$ be three triangulated categories, $\cA$ a thick
subcategory of $\cB$, $\cB$ a thick subcategory of $\cC$.
Then $\cA$ is a thick subcategory of $\cC$, $\cB/\cA$ is a
thick subcategory of $\cC/\cA$. The canonical functor
$\cC/\cB \to (\cC/\cA)/(\cB/\cA)$ is an isomorphism. 
\end{corollary}










\subsection{Properties of the quotient functor}\label{VIII:2-5}

Let $\cA$ be a triangulated category, $\cB$ a thick subcategory,
$S$ the corresponding multiplicative system, $Q:\cA\to \cA/\cB$ the canonical
quotient functor. 





\begin{proposition}\label{VIII:2-5-1}
Given objects $X$ and $Y$ of $\cA$, the following assertions are
equivalent:
\begin{enumerate}[(a)]
  \item Every morphism $f:X\to Y$, such that there exists a morphism $s:Z\to X$, in
    $S$, satisfying $f\circ s=0$, is a zero morphism.
  \item Every morphism $f:X\to Y$, such that there exists a morphism $s:Y\to Z$,
    $s\in S$, satisfying $s\circ f=0$, is a zero morphism.
  \item Every morphism $f:X\to Y$ factoring through an object of $\cB$ is
    zero.
  \item The canonical map: 
    \[
      \hom_\cA(X,Y) \to \hom_{\cA/\cB}(Q(X),Q(Y))
    \]
    is injective.
\end{enumerate}
\end{proposition}





\begin{proposition}\label{VIII:2-5-2}
Let $X$ and $Y$ be two objects of $\cA$. The following assertions are
equivalent:
\begin{enumerate}[(a)]
  \item Every diagram ($s\in S$): 
    \[\xymatrix{
      & Z \ar[dl]_-s \ar[dr]^-f \\
      X 
        & & Y 
    }\]
    completes to a diagram: $(s,t\in S$): 
    \[\xymatrix{
      & Z' \ar[d]_-t \\
      & Z \ar[dl]_-s \ar[dr]^-f \\
      X 
        & & Y 
    }\]
    where $f\circ t$ factors through $(X,s\circ t)$.
  \item Assertion obtained by reversing the direction of arrows in (a) and
    interchanging $X$ and $Y$.
  \item Every diagram: 
    \[\xymatrix{
      X \ar[r]^-f 
        & N \ar[r]^-g 
        & Y 
    }\]
    where $g\circ f=0$ and where $N$ is an object of $\cB$, completes to a
    diagram: 
    \[\xymatrix{
      & N' \ar[d]^-i \\
      X \ar[ur]^-h \ar[r]^-f 
        & N \ar[r]^-g 
        & Y 
    }\]
    where $f=i\circ h$ and where $g\circ i=0$.
  \item Assertion obtained by changing the direction of arrows in (c) and
    interchanging $X$ and $Y$.
  \item The canonical map: 
    \[\xymatrix{
      \hom_\cA(X,Y) \ar[r]^-Q 
        & \hom_{\cA/\cB}(Q(X),Q(Y)) 
    }\]
    is surjective.
\end{enumerate}
\end{proposition}





\begin{proposition}\label{VIII:2-5-3}
Let $X$ be an object of $\cA$. The following assertions are equivalent:
\begin{enumerate}[(a)]
  \item For every object $Y$ of $\cA$, the canonical map
    \[\xymatrix{
      \hom_\cA(X,Y) \ar[r]^-Q
        & \hom_{\cA/\cB}(Q(X),Q(Y))
    }\]
    is an isomorphism.
  \item Every morphism from $X$ to an object of $\cB$ is zero.
\end{enumerate}
Similarly, the two following assertions are equivalent:
\begin{enumerate}[(a)']
  \item For every object $Y$ of $\cA$, the canonical map:
    \[\xymatrix{
      \hom_\cA(X,Y) \ar[r]^-Q
        & \hom_{\cA/\cB}(Q(X),Q(Y))
    }\]
    is an isomorphism.
  \item Every morphism from an object of $\cB$ to $X$ is zero.
\end{enumerate}
\end{proposition}





\begin{definition}\label{VIII:2-5-4}
Every object satisfying the equivalent properties (a) and (b) of the
preceding proposition, will be called: \emph{left free object over
$\cA/\cB$} or \emph{left $Q$-free} object. Similarly, every
object satisfying the equivalent properties (a)' and (b)' will be called
\emph{right free object over $\cA/\cB$} or
\emph{right $Q$-free} object. It is defined up to isomorphism by the
knowledge of $Q(X)\in\operatorname{Ob}{\cA/\cB}$.
\end{definition}





\subsection{Functors adjoint to the quotient functor}\label{VIII:2-6}





\begin{definition}\label{VIII:2-6-1}
Two triangulated subcategories $\cN$ and $\cN'$ of a
triangulated category $\cA$ are said to be \emph{orthogonal} if for every object $X$ of
$\cN$ and every object $Y$ of $\cN'$, we have:
\[
  \hom_\cA(X,Y) = 0
\]
$\cN'$ will then be said to be right orthogonal to $\cN$ and $\cN$ left
orthogonal to $\cN'$.
\end{definition}





\begin{proposition}\label{VIII:2-6-2}
Let $\cN$ be a triangulated subcategory of a triangulated category
$\cA$. The full subcategory $\cN^\bot$ (resp. $^\bot\cN$) generated by the
objects $X$ of $\cA$ such that for every object $Y$ of $\cN$ we have
$\hom_\cA(Y,X)=0$ (resp. $\hom_\cA(X,Y)=0$), is a thick subcategory
of $\cA$. The category $\cN^\bot$ is called by abuse of language,
the right orthogonal of $\cN$.
\end{proposition}

This proposition allows us to restate proposition \ref{VIII:2-5-3}. 





\begin{proposition}\label{VIII:2-6-3}
Let $\cB\subset\cA$, be a thick subcategory of a
triangulated category $\cA$. The full category generated by the right free
(resp. left free) objects over $\cA/\cB$, is none other than the right
(resp. left) orthogonal of $\cB$. 
\end{proposition}





\begin{proposition}\label{VIII:2-6-4}
Let $\cA$ be a triangulated category, $\cB$ a thick
subcategory of $\cA$, $\cB^\bot$ the right orthogonal of $\cB$, $S(\cB)$,
$s(\cB^\bot)$ the corresponding multiplicative systems, $Q$ and $Q^\bot$ the
canonical quotient functors. Consider for an object $x$ of
$\cA$, the following properties:
\begin{enumerate}[(i)]
  \item $X$ maps by a morphism of $S(\cB)$ to an object of $\cB^\bot$.
  \item $X$ receives by a morphism of $S(\cB^\bot)$ an object of $\cB$.
  \item The category $S_X(\cB)$ of arrows of $S(\cB)$ with source $X$,
    admits a final object.
  \item The category $S^X(\cB^\bot)$ of arrows of $S(\cB^\bot)$ with target
    $X$, admits an initial object.
  \item The category $\cB/X$ of objects of $\cB$ over $X$, admits a
    final object.
  \item The category $\cB^\bot/X$ of objects of $\cB^\bot$ under $X$
    admits an initial object.
\end{enumerate}
We then have 
\[\xymatrix@=0.5cm{
  & & & & (iv) \\
  (i) \ar@{<=>}[r] 
    & (ii) \ar@{<=>}[r] 
    & (iii) \ar@{<=>}[r] 
    & (v) \ar@{=>}[ur] \ar@{=>}[dr] \\
  & & & & (vi) 
}\]

If moreover $^\bot(\cB^\bot)=\cB$ then all the properties are
equivalent.
\end{proposition}

Let $\cB$ be a thick subcategory of a triangulated category
$\cA$, the inclusion functor $i$ of $\cB$ into $\cA$ and $Q$ the
quotient functor from $\cA$ to $\cA/\cB$, $\cB^\bot$ the right
orthogonal of $\cB$, the corresponding functors $i^\bot$ and $Q^\bot$. One deduces
immediately from proposition \ref{VIII:2-6-4} the following propositions: 





\begin{proposition}\label{VIII:2-6-5}
The two following properties are equivalent:
\begin{enumerate}[(i)]
  \item The functor $i$ admits a right adjoint.
  \item The functor $Q$ admits a right adjoint.
\end{enumerate}
\end{proposition}

The proposition remains true when one replaces the word right by the word
left. 





\begin{proposition}\label{VIII:2-6-6}
The two following properties are equivalent:
\begin{enumerate}[(i)]
  \item The functor $i$ admits a right adjoint.
  \item The functor $i^\bot$ admits a left adjoint. The left
    orthogonal of the category $\cB^\bot$ equals $\cB$. 
\end{enumerate}
\end{proposition}





\begin{proposition}\label{VIII:2-6-7}
Suppose the properties of propositions \ref{VIII:2-6-5} and
\ref{VIII:2-6-6} hold. Let $i^\ast$ and $Q^\ast$ be the right adjoints of $i$ and
$Q$. Likewise let $^\ast i^\bot$ and $^\ast Q^\bot$ be the left adjoints
of $i^\bot$ and $Q^\bot$. All these functors are exact. Moreover:

There exists a functorial isomorphism
$Q^\ast\circ Q\iso i^\bot\circ ^\ast i^\bot$ such that the diagram below
is commutative: 
\[\xymatrix{
  & \operatorname{id} \ar[dl] \ar[dr] \\
  Q^\ast \circ Q \ar[rr]^\sim 
    & & i^\bot\circ ^\ast i^\bot 
}\]
(the oblique arrows are defined by the adjunction properties).

There likewise exists an isomorphism
$^\ast Q^\bot \circ Q^\bot \iso i\circ i^\ast$ such that the following diagram is
commutative: 
\[\xymatrix{
  ^\ast Q^\bot\circ Q^\bot \ar[rr]^-\sim \ar[dr] 
    & & i\circ i^\ast \ar[dl] \\
  & \operatorname{id}{}
}\]

Finally, there exists a functorial morphism of degree $1$:
$\partial:i^\bot\circ ^\ast i^\bot \to i\circ i^\ast$ such that the following
triangle is distinguished $(\deg\partial=1$): 
\[\xymatrix{
  i^\bot \circ ^\ast i^\bot \ar[rr]^-\partial 
    & & i\circ i^\ast \ar[dl] \\
  & \operatorname{id}{} \ar[ul] 
}\]

Such a morphism $\delta$ is unique and satisfies
$\delta(T X_0) = - T \delta(X)$ ($X$ object $q\circ q$, $T$ translation).
\end{proposition}















\section{Derived categories of an abelian category}\label{VIII:3}





\subsection{Derived categories}\label{VIII:3-1}

We use the notation of \ref{VIII:1-2-3}. We shall moreover use the
following notation:

\subsubsection{Notation}\label{VIII:3-1-1}

Let $\cA$ be an abelian category. We set: 
\begin{align*}
  \eD(\cA) &= \eK^{\infty,\infty}(\cA)/\eK^{\infty,\varnothing}(\cA) \\
  \eD^b(\cA) &= \eK^{b,b}(\cA)/\eK^{b,\varnothing}(\cA) \\
  \eD^+(\cA) &= \eK^{+,+}(\cA) / \eK^{+,\varnothing}(\cA) \\
  \eD^-(\cA) &= \eK^{-,-}(\cA) / \eK^{-,\varnothing}(\cA) \text{.}
\end{align*}





\begin{proposition}\label{VIII:3-1-2}
The natural functors appearing in the following diagram:
\[\xymatrix@=0.5cm{
  & \eD(\cA) \\
  \eD^+(\cA) \ar[ur]
    & & \eD^-(\cA) \ar[ul] \\
  & \eD^b(\cA) \ar[ul] \ar[ur]
}\]
are fully faithful.

The natural functors: 
\begin{align*}
  \eD^b(\cA) &= \eK^{b,b}(\cA)/\eK^{b,\varnothing}(\cA) \to \eK^{+,b}(\cA)/\eK^{+,\varnothing}(\cA) \\
  \eD^b(\cA) &= \eK^{b,b}(\cA)/\eK^{b,\varnothing}(\cA) \to \eK^{-,b}(\cA) / \eK^{-,\varnothing}(\cA) 
\end{align*}
are equivalences of categories.

The functors in the diagram: 
\[\xymatrix@=0.5cm{
  & \eD(\cA) \\
  \eD^+(\cA) \ar[ur] 
    & & \eD^-(\cA) \ar[ul] \\
  & \eD^b(\cA) \ar[ul] \ar[ur]
}\]
being injective on objects, realise the categories $\eD^b(\cA)$,
$\eD^+(\cA)$, $\eD^-(\cA)$ as full subcategories of $\eD(\cA)$.
\end{proposition}

The proof of this proposition is found in theorem \ref{VIII:2-4-2}. 





\begin{definition}\label{VIII:3-1-3}
The category $\eD(\cA)$ will be called the \emph{derived category} of the
category $\cA$. The objects of $\eD(\cA)$ isomorphic to objects of $\eD^b(\cA)$
will be called the \emph{bounded} objects of $\eD(\cA)$. The objects of
$\eD(\cA)$ isomorphic to objects of $\eD^+(\cA)$ will be called
\emph{left bounded} objects, the objects of $\eD(\cA)$ isomorphic to objects of
$\eD^-(\cA)$, will be called \emph{right bounded} objects.
\end{definition}

We denote by $D$ the canonical functor:
\[
  D : \eC(\cA) \to \eD(\cA) \text{.}
\]
The ``restriction'' of $D$ to the full subcategory $\cA$ of $\eC(\cA)$
will still be denoted $D$. The functor $D$ restricted to $\cA$ is fully
faithful. The functor $D$ factors uniquely through $\eK(\cA)$.
Let $X$ and $Y$ be two objects of $\cA$.

One checks without difficulty that for every $m\geqslant 0$, the groups:
\[
  \hom_{\eD(\cA)}(D(X),T^{-m} D(Y)) \qquad \text{($T$ is the translation functor)}
\]
vanish. The cohomological dimension of $\cA$ will be the smallest integer
$n$ such that for every $m>n$ we have:
\[
  \hom_{\eD(\cA)}(D(X),T^m D(Y)) = 0 \qquad\text{for every pair $X,Y$ of objects of $\cA$.}
\]
If there is no such integer, we say that the cohomological dimension of $\cA$
is infinite.

Denote by $\eI^+(\cA)$ (resp. $\eI(\cA)$) the full triangulated subcategory
of $\eK^+(\cA)$ (resp. of $\eK(\cA)$) defined by the complexes whose
objects are injective in every degree. 





\begin{proposition}\label{VIII:3-1-4}
Suppose that the category $\cA$ has enough injectives. The
natural functor:
\[
  Q^+:\eK^+(\cA) \to \eD^+(\cA)
\]
induces an equivalence of categories:
\[
  \eI^+(\cA) \to \eD^+(\cA)  \text{.}
\]
The quasi-inverse functor is a right adjoint of $Q^+$.
\end{proposition}

If moreover $\cA$ is of finite cohomological dimension, the preceding
assertion remains true when the exponents $+$ are dropped.

There is an analogous statement for projectives.

The proof of proposition \ref{VIII:3-1-4} relies on
\ref{VIII:2-6-4}.

The triangulated structure of $\eD(\cA)$, is made precise by the following
proposition.

Let $\eE(\cA)$ be the category of three-term exact sequences of
$\eC(\cA)$. The category $\mathsf{SSS}(\cA)$ of semi-split exact
sequences (\S\ref{VIII:1-2-4}) is a full subcategory of
$\eE(\cA)$. We have defined (loc.\ cit.) a functor:
\[
  \rho:\mathsf{SSS}(\cA) \to \mathsf{TrK}(\cA)
\]
where $\mathsf{TrK}(\cA)$ denotes the category of distinguished triangles
of $\eK(\cA)$. One deduces a functor
\[
  \sigma : \mathsf{SSS}(\cA) \to \mathsf{TrD}(\cA)
\]
where $\mathsf{TrD}(\cA)$ denotes the category of distinguished triangles
of $\eD(\cA)$. 





\begin{proposition}\label{VIII:3-1-5}
There exists one and only one functor:
\[
  \Pi:\eE(\cA) \to \mathsf{TrD}(\cA)
\]
of the form ($\deg \delta(S)=1$): 
\[
(S=0 \to X^\bullet \to Y^\bullet \to Z^\bullet \to 0) 
\mapsto 
\xymatrix{
  & D(Z^\bullet) \ar[dl]_-{\delta(S)} \\
  D(X^\bullet) \ar[rr]^-{D(u)} 
    & & D(Y^\bullet) \ar[ul]_-{D(v)} 
}
\]
whose restriction to $\mathsf{SSS}(\cA)$ is $\sigma$. \emph{This functor
is essentially surjective.}
\end{proposition}





\begin{proposition}\label{VIII:3-1-6}
The canonical cohomological functor:
\[
  \h^0:\eK(\cA) \to \cA
\]
factors uniquely through a functor which we still denote
$\h^0$
\[
  \h^0:\eD(\cA) \to \cA \text{.}
\]
On $\eD(\cA)$ the functor $\h^0$ and its translates $\h^i$ form a
conservative system, i.e. a morphism $f:X^\bullet \to Y^\bullet$ is an
isomorphism if and only if for every integer $i$
\[
  \h^i(f):\h^i(X) \to \h^i(Y)
\]
is an isomorphism.
\end{proposition}





\begin{definition}\label{VIII:3-1-7}
A morphism $f:X^\bullet \to Y^\bullet$ of $\eC(\cA)$ (resp. of $\eK(\cA)$)
is called a \emph{quasi-isomorphism} when for every integer $i$
\[
  \h^i(f):\h^i(X) \to \h^i(Y)
\]
is an isomorphism.
\end{definition}

By the preceding proposition quasi-isomorphisms are the
morphisms which become isomorphisms in $\eD(\cA)$. 





\begin{definition}\label{VIII:3-1-8}
Let $\sigma$ be a set of objects of $\cA$, $X^\bullet$ an object of
$\eC(\cA)$ (resp. $\eK(\cA)$). A right (resp. left) resolution of type
$\sigma$ of $X^\bullet$, is a quasi-isomorphism $X^\bullet \to V^\bullet$
(resp. $V^\bullet \to X^\bullet$) where $V^\bullet$ is a complex all of
whose objects lie in $\sigma$. We say \emph{injective resolution}
(resp. projective) instead of right (resp. left) resolution of
injective (resp. projective) type.
\end{definition}

In general, we shall allow ourselves to drop the words ``right'' or
``left'' when no ambiguity results. 





\begin{proposition}\label{VIII:3-1-9}
\begin{enumerate}
  \item Suppose that every object of $\cA$ embeds into an object of $\sigma$
    (resp. is a quotient of an object of $\sigma$). Then every object of
    $\eC^+(\cA)$ (resp. of $\eC^-(\cA)$) admits a right
    (resp. left) resolution of type $\sigma$ in $\eC^+(\cA)$ (resp. in
    $\eC^-(\cA)$).
  \item Suppose moreover that there exists an integer $n$ such that for every exact
    sequence: 
    \begin{align*}
      Y^0 &\to Y^1 \to \cdots \to Y^{n-1} \to Y^n \to 0 \\
      \text{(resp. } 0 &\to Y^n \to Y^{n-1} \to \cdots \to Y^1 \to Y^0 \text{)} 
    \end{align*}
    of objects of $\cA$ where for every $i$, $0\leqslant i\leqslant n-1$, $Y^i$ is
    an object of $\sigma$, the object $Y^n$ is an object of $\sigma$. Then every
    object of $\eC(\cA)$ admits a right (resp. left) resolution of type
    $\sigma$. 
\end{enumerate}
\end{proposition}










\subsection{Study of $\ext$}\label{VIII:3-2}

Let $X^\bullet$ and $Y^\bullet$ be two objects of $\eC(\cA)$ (resp. $\eK(\cA)$)
and $D:\eC(\cA) \to \eD(\cA)$ (resp. $Q:\eK(\cA) \to \eD(\cA)$) the canonical
functor. 





\begin{definition}\label{VIII:3-2-1}
We call \emph{$i$-th hyper-ext} and denote: 
\[
  (X^\bullet,Y^\bullet)\mapsto \ext^i(X^\bullet,Y^\bullet) 
\]
the functor: $\hom_{\eD(\cA)}(D(X^\bullet),T^i D(Y^\bullet))$ (resp.
$\hom_{\eD(\cA)}(Q(X^\bullet),T^i Q(Y^\bullet))$). 
\end{definition}





\begin{proposition}\label{VIII:3-2-2}
Let $0 \to X^\bullet \to Y^\bullet \to Z^\bullet \to 0$ be an exact sequence of
$\eC(\cA)$. Let $V^\bullet$ be an object of $\eC(\cA)$. We have the unbounded
exact sequences: 
\[\xymatrix@=0.5cm{
  \cdots \ar[r] 
    & \ext^i(V^\bullet,X^\bullet) \ar[r] 
    & \ext^i(V^\bullet,Y^\bullet) \ar[r] 
    & \ext^i(V^\bullet,Z^\bullet) \ar[r]^-\delta 
    & \ext^{i+1}(V^\bullet,X^\bullet) \ar[r] 
    & \cdots \\
  \cdots \ar[r] 
    & \ext^i(Z^\bullet,V^\bullet) \ar[r] 
    & \ext^i(Y^\bullet,V^\bullet) \ar[r] 
    & \ext^i(X^\bullet,V^\bullet) \ar[r]^-\delta 
    & \ext^{i+1}(Z^\bullet,V^\bullet) \ar[r] 
    & \cdots
}\]
\end{proposition}

Let $X^\bullet$ and $Y^\bullet$ be two objects of $\eK(\cA)$. We denote by
$\mathsf{Qis}/X^\bullet$ (resp. $\mathsf{Qis}^+/X^\bullet$,
$\mathsf{Qis}^-/X^\bullet$, $\mathsf{Qis}^b/X^\bullet$) the category of
quasi-isomorphisms with target $X^\bullet$ and source in $\eK(\cA)$ (resp.
$\eK^+(\cA)$, $\eK^-(\cA)$, $\eK^b(\cA)$). Likewise we define the
categories $Y^\bullet/\mathsf{Qis}$, $Y^\bullet/\mathsf{Qis}^+$,\ldots
(quasi-isomorphisms with source $Y^\bullet$).

We propose to summarise in the following proposition the various
equivalent definitions of $\ext^i$. 





\begin{proposition}\label{VIII:3-2-3}
1. There exist isomorphisms of functors: 
\[\xymatrix{
  \ext^0(X^\bullet,Y^\bullet) = \hom_{\eD(\cA)}(Q(X^\bullet),Q(Y^\bullet)) \ar[r]^-\sim 
    & \varinjlim_{(\mathsf{Qis}/X^\bullet)^\circ} \hom_{\eK(\cA)}(-,Y^\bullet) \ar[r]^-\sim 
    & \\
  \varinjlim_{Y^\bullet/\mathsf{Qis}} \hom_{\eK(\cA)}(X^\bullet,-) \ar[r]^-\sim 
    & \varinjlim_{(\mathsf{Qis}/X^\bullet)^\circ \times Y^\bullet/\mathsf{Qis}} \hom_{\eK(\cA)}(-,-) \text{.}
}\]
If $X^\bullet$ is an object of $\eK^-(\cA)$: 
\[
  \ext^0(X^\bullet,Y^\bullet) \iso \varinjlim_{(\mathsf{Qis}^-/X^\bullet)^\circ} \hom_{\eK(\cA)}(-,Y^\bullet) 
\]
If $Y^\bullet$ is an object of $\eK^+(\cA)$: 
\[
  \ext^0(X^\bullet,Y^\bullet) \iso \varinjlim_{Y^\bullet/\mathsf{Qis}} \hom_{\eK(\cA)}(X^\bullet,-) 
\]
If $X^\bullet$ is an object of $\eK^-(\cA)$ and if $Y^\bullet$ is an
object of $\eK^b(\cA)$: 
\[
  \ext^0(X^\bullet,Y^\bullet) \iso \varinjlim_{Y^\bullet/\mathsf{Qis}^b} \hom_{\eK(\cA)}(X^\bullet,-) 
\]
If $X^\bullet$ is an object of $\eK^b(\cA)$ and $Y^\bullet$ an object of
$\eK^+(\cA)$: 
\[
  \ext^0(X^\bullet,Y^\bullet) \iso \varinjlim_{(\mathsf{Qis}^b/X^\bullet)^\circ} \hom_{\eK(\cA)}(-,Y^\bullet) 
\]
If the category $\cA$ has enough injectives and if
$Y^\bullet$ is an object of $\eK^+(\cA)$, $Y^\bullet$ admits an
injective resolution and only one (up to isomorphism in
$\eK^+(\cA)$): 
\[
  Y^\bullet \to I(Y^\bullet) 
\]
and we have:
\[
  \ext^0(X^\bullet,Y^\bullet) \iso \hom_{\eK(\cA)}(X^\bullet,I(Y^\bullet)) \text{.}
\]
We likewise have an analogous statement for projectives.
If the category $\cA$ admits enough injectives and if $\cA$ is
of \emph{finite cohomological dimension}, every object $Y^\bullet$ of
$\eK(\cA)$ admits one and only one injective resolution:
\[
  Y^\bullet \to I(Y^\bullet)
\]
and we have:
\[
  \ext^0(X^\bullet,Y^\bullet) \iso \hom_{\eK(\cA)}(X^\bullet,I(Y^\bullet))
\]
Analogous statement for projectives.
\end{proposition}

\paragraph{Remark:}
Unfolding the isomorphism of proposition \ref{VIII:3-2-3}(3), one
recovers Yoneda's definition. 










\section{Derived functors}\label{VIII:4}










\subsection{Definition of derived functors}\label{VIII:4-1}





\begin{definition}\label{VIII:4-1-1}
Let $\cC$ and $\cC'$ be two graded categories (we denote by $T$ the
translation functor of $\cC$ and of $\cC'$), $F$ and $G$ two graded
functors from $\cC$ to $\cC'$. A \emph{morphism of graded functors} is
a morphism of functors:
\[
  u:F\to G
\]
with the following property:

For every object $X$ of $\cC$ the following diagram commutes: 
\[\xymatrix{
  F(T X) \ar[r]^-{u(T X)} 
    & G(T X) \\
  T F(X) \ar[r]^-{T u(X)} \ar[u]_-\wr 
    & T G(X) \ar[u]_-\wr 
}\]
\end{definition}

Let $\cC$ and $\cC'$ be two triangulated categories. We denote by
$\mathsf{Fex}(\cC,\cC')$ the category of exact functors from $\cC$ to
$\cC'$, the morphisms between two functors being the morphisms of graded
functors.

Let $\cA$ and $\cB$ be two abelian categories and
$\Phi:\eK^\ast(\cA)\to \eK^{\ast'}(\cB)$ an exact functor ($\ast$ and $\ast'$
denote one of the signs $+$, $-$, $b$, or $v$ ``empty''). The canonical
functor $Q:\eK^\ast(\cA)\to \eD^\ast(\cA)$ gives us, by composition, a
functor
\[\xymatrix{
  \mathsf{Fex}(\eD^\ast(\cA),\eD^{\ast'}(\cB)) \ar[r]
    & \mathsf{Fex}(\eK^\ast(\cA),\eD^{\ast'}(\cB))
}\]
whence (also denoting by $Q'$ the canonical functor
$\eK^{\ast'}(\cB) \to \eD^{\ast'}(\cB)$) a functor:
$\%$ (resp. $\%'$):
$\mathsf{Fex}(\eD^\ast(\cA),\eD^{\ast'}(\cB)) \to \mathsf{Ab}$: 
\begin{align*}
  \Psi  &\mapsto \hom(Q'\circ \Phi,\Psi\circ Q) \\
  (\text{resp. }\Psi &\mapsto \hom(\Psi\circ Q,Q'\circ \Phi)\text{ ).} 
\end{align*}





\begin{definition}\label{VIII:4-1-2}
We say that $\Phi$ admits a \emph{total right} (resp.
left) \emph{derived functor} if the functor $\%$ (resp. $\%'$) is representable. An
object representing the functor $\%$ (resp. $\%'$) will be called total right
(resp. left) derived functor of $\Phi$ and will be denoted
$\eR \Phi$ (resp. $\eL \Phi$).\footnote{A more manageable definition (and in
practice equivalent) is given in \cite[XVII 1.2]{sga4}.}
\end{definition}





\subsubsection{Notation}\label{VIII:4-1-3}

Let $\cA$ and $\cB$ be two abelian categories and $f:\cA\to \cB$ an
additive functor. The \emph{total right derived functor} of the functor
$\eK^+(f):\eK^+(\cA) \to \eK^+(\cB)$ will be denoted $\eR^+ f$. One defines
likewise $\eR^- f$, $\eR f$, $\eL^+ f$, $\eL^- f$, $\eL f$.

Likewise let $F:\eK^\ast(\cA) \to \cB$, be a cohomological functor. The
canonical quotient functor $Q:\eK^\ast(\cA) \to \eD^\ast(\cA)$
defines by composition a functor: 
\[\xymatrix{
  \mathsf{Foco}(\eD^\ast(\cA),\cB) \ar[r] 
    & \mathsf{Foco}(\eK^\ast(\cA),\cB) 
}\]
($\mathsf{Foco}(-,-)$ denotes the category of cohomological functors)
whence a functor $\%$ (resp. $\%'$)
$\mathsf{Foco}(\eD^\ast(\cA),\cB) \to \mathsf{Ab}$, $G\mapsto \hom(F,G\circ Q)$
(resp. $G\mapsto \hom(G\circ Q,F)$). 





\begin{definition}\label{VIII:4-1-4}
We say that the functor $F$ admits a \emph{cohomological derived functor}
on the right (resp. on the left) if the functor $\%$ (resp. $\%'$) is
representable. An object representing the functor $\%$ (resp. $\%'$) will be
called right (resp. left) cohomological derived functor and
will be denoted $\eR F$ (resp. $\eL F$). 
\end{definition}





\subsubsection{Notation}\label{VIII:4-1-5}

Let $\cA$ and $\cB$ be two abelian categories and $f:\cA\to \cB$ an
additive functor. The right cohomological derived functor of the functor:
\[
  \h^0{}\circ \eK^+ f:\eK^+(\cA) \to \cB
\]
will be denoted $\eR^+ f$. When no confusion results the functor
$\eR^+ f\circ T^i$ ($T$ is the translation functor) will be denoted:
$\eR^i f$. The right cohomological derived functor of the functor
$\h{}\circ \eK f$ will be denoted $\eR f$. When no confusion results
the functor $\eR f\circ T^i$ will be denoted $\eR^i f$.

One defines likewise the functors $\eL^- f$, $\eL f$, $\eL^i f$. 





\subsubsection{Remarks}\label{VIII:4-1-6}

These definitions contain, as we shall see in the next
section, the classical definition of derived functors when the
categories have enough injectives or projectives and the complexes
are suitably bounded. They are however more general. But
without further hypotheses they are hard to handle. For example, let
$f:\cA\to \cB$ be an additive functor between two abelian categories.
Suppose that the functors $\eR^+ f$ and $\eR f$ exist. I do not know how to
prove that $\eR f$ induces $\eR^+ f$ on $\eD^+(\cA)$. Suppose
moreover that the functor $\eR^+ f$ (cohomological derived) exists. I do not know
how to prove that the functor $\h^0{}\circ \eR^+ f$ is isomorphic to the functor
$\eR^+ f$. In practice however these properties will hold. 










\subsection{Existence of derived functors}\label{VIII:4-2}





\begin{proposition}\label{VIII:4-2-1}
The hypotheses are those of definition \ref{VIII:4-1-4}. Suppose
moreover that $\cA$ is a $\sU$-category, where $\sU$ is a universe such
that the set $\dZ$ of rational integers is an element of $\sU$, such
that the set of objects of $\cA$ is an element of $\sU$. Suppose
moreover that the category $\cB$ is the category of $\sU$-abelian
groups or more generally the category of sheaves of
$\sU$-abelian groups over an arbitrary $\sU$-site. (A $\sU$-site is a
$\sU$-category, whose set of objects is an element of $\sU$,
with a topology.) The functor $F$ admits a right
cohomological derived functor. This functor is computed as follows:

Let $X$ be an object of $\eD^\ast(\cA)$, $\underline X$ the object of
$\eK^\ast(\cA)$ above $X$. We then have a functorial isomorphism:
\[
  \eR F(X) = \varinjlim_{\underline X/\mathsf{Qis}^\ast} F(-) \text{.}
\]
\end{proposition}

In particular let $f:\cA\to \cB$ be an additive functor. The functors
$\eR f$ and $\eR^+ f$ exist and are computed as indicated
above. The functor $\eR f$ induces $\eR^+ f$ on $\eD^+(\cA)$.
As no confusion can result, we
use the notation $\eR^i f$. 





\subsubsection{Remarks}\label{VIII:4-2-2}

One can state proposition \ref{VIII:4-2-1} in a more
general form in the setting of triangulated categories. The hypotheses
to impose on the category $\cB$, to ensure the validity of
proposition \ref{VIII:4-2-1} are hypotheses of existence and exactness
of inductive limits over pseudo-filtered $\sU$-categories, whose
sets of objects are elements of $\sU$. Proposition
\ref{VIII:4-2-1} introduces a dissymmetry between right derived functors
and left derived functors at least in practice. 





\begin{theorem}\label{VIII:4-2-3} % labelled 4.2.2 in original
The data are those of definition \ref{VIII:4-1-2}. We are further given a full
triangulated subcategory of $\eK^\ast(\cA)$: $\cC$, and we
denote by $\cN$ the full triangulated subcategory of objects of $\cC$
which are acyclic. Suppose that
\begin{enumerate}[(1)]
  \item Every object of $\cN$ is sent by $\Phi$ to an acyclic object
    of $\eK^{\ast'}(\cB)$.
  \item Every object $X$ of $\eK^\ast(\cA)$ is the source (resp. target) of a
    quasi-isomorphism whose target (resp. source) is an object of $\cC$.
\end{enumerate}
\begin{enumerate}[a)]
  \item Then $\Phi$ admits a total right (resp.
    left) derived functor.
  \item The functor $\eR \Phi$ (resp. $\eL \Phi$) is obtained as
    follows:

    The restriction of the functor
    $Q'\circ \Phi:\eK^\ast(\cA) \to \eD^{\ast'}(\cB)$ to the category $\cC$
    vanishes on the objects of $\cN$, hence factors through $\cC/\cN$. We
    thus obtain a functor $\Phi':\cC/\cN\to \eD^{\ast'}(\cB)$. But the
    natural functor: $\cC/\cN\to \eD^\ast(\cA)$ is an equivalence of
    categories (theorem \ref{VIII:2-4-2}). Whence, composing with the
    quasi-inverse functor, the functor $\eR\Phi$ (resp. $\eL \Phi$).
  \item Let $X$ be an object of $\eK^\ast(\cA)$. Let $Y$ be an object of $\cC$ and
    $X\to Y$ (resp. $Y\to X$) a quasi-isomorphism. The object
    $\eR\Phi\circ Q(X)$ (resp. $\eL \Phi\circ Q(X)$) is isomorphic to
    $Q\circ \Phi(Y)$.
  \item The functor $X\mapsto \eR\Phi\circ Q(X)$ (resp.
    $X\mapsto \eL \Phi\circ Q(X)$) is isomorphic to the functor:
    $X\mapsto \varinjlim_{X/\mathsf{Qis}^\ast} Q\circ\Phi(-)$ (resp.
    $X\mapsto \varinjlim_{\mathsf{Qis}^\ast/X} Q\circ \Phi(-)$).
  \item The functor $\h^0{}\circ \Phi:\eK^\ast(\cA)\to \cB$, admits a
    right (resp. left) cohomological derived functor which is
    none other than the functor $\h^0{}\circ\eR\Phi$ (resp. $\h^0{}\circ \eL\Phi$). 
\end{enumerate}
\end{theorem}





\begin{corollary}\label{VIII:4-2-3-1} % originally corollary 1
Suppose that $\ast=+$ (resp. $\ast=-$) in the hypotheses of the preceding
theorem, and that $\cA$ has enough injectives (resp.
projectives). Then theorem \ref{VIII:4-2-3} applies taking for
$\cC$ the category of injective (resp. projective) complexes. Suppose moreover
that $\cA$ is of finite cohomological dimension, then the theorem
applies, with no restriction on $\ast$, taking for $\cC$ the category
of injective (resp. projective) complexes. 
\end{corollary}





\begin{definition}\label{VIII:4-2-5} % originally 2.4
Let $\cA$ and $\cB$ be two abelian categories and $f:\cA\to \cB$ an
additive functor.

We say that $f$ is of \emph{finite cohomological dimension} on the right (resp.
on the left) if $\eR^+ f$ (resp. $\eL^- f$) exists and if there exists an integer
$m\geqslant 0$ such that for every integer $n>m$, $\eR^n f$ (resp.
$\eL^{-n} f$) vanishes on objects coming from $\cA$ via the functor $D$.
The right (resp. left) \emph{cohomological dimension} of $f$ is
then the smallest integer $m$ with the above property. If
$\eR^+ f$ (resp. $\eL^- f$) exists and $f$ is not of finite
cohomological dimension, we say that $f$ is of infinite cohomological dimension.

An object $X$ of $\cA$ will be said to be \emph{$f$-acyclic} on the right (resp. on the
left) if for every integer $n>0$, $\eR^n f(D(X)) = 0$ (resp.
$\eL^{-n} f(D(X))=0$). 
\end{definition}





\begin{corollary}\label{VIII:4-2-3-2} % originally corollary 2
The data are those of definition \ref{VIII:4-2-5}. Suppose that there
exists a set $\sigma$ of objects of $\cA$ stable under direct sum
with the following properties:
\begin{enumerate}
  \item Every acyclic complex $\in \operatorname{Ob}{\eK^+(\cA)}$ (resp.
    $\in\operatorname{Ob}{\eK^-(\cA)}$) of objects of $\sigma$, is sent
    by $f$ to an acyclic complex.
  \item Every object $X$ of $\cA$ embeds into (resp. is a quotient of) an object
    of $\sigma$. The hypotheses of theorem \ref{VIII:2-4-3} hold
    taking $\ast=+$ (resp. $\ast=-$) and taking for
    category $\cC$, the category of complexes whose objects in every
    degree are elements of $\sigma$. The elements of $\sigma$
    are right (resp. left) $f$-acyclic objects.
\end{enumerate}

If moreover $f$ is of finite cohomological dimension on the right (resp. on the
left), the set of right (resp. left) $f$-acyclic objects
has in addition to properties (1) and (2) above, property
(2) of proposition \ref{VIII:3-1-9}. Theorem \ref{VIII:4-2-3}
then applies to the functor $\eK f$ taking for category $\cC$ the
category of complexes whose objects are right
(resp. left) $f$-acyclic in every degree. The functor $\eR f$ (resp. $\eL f$) exists
and induces on $\eD^+(\cA)$ the functor $\eR^+ f$ (resp. $\eL^+ f$) and on
$\eD^-(\cA)$ the functor $\eR^- f$ (resp. $\eL^- f$). 
\end{corollary}










\subsection{Product. Composition}\label{VIII:4-3}

Let $\cA$ be an abelian category, $F:\eD(\cA) \to \mathsf{Ab}$ a
cohomological functor. Let $X^\bullet$ and $Y^\bullet$ be two objects of
$\eC(\cA)$. The functor $F$ defines a map:
\[
  \ext^0(X^\bullet,Y^\bullet) \to \hom_{\mathsf{Ab}}(F(D(X^\bullet)),F(D(Y^\bullet)))
\]
whence a pairing:
$F(D(X^\bullet))\times \ext^0(X^\bullet,Y^\bullet) \to F(D(Y^\bullet))$.

This pairing gives, applying translations, pairings:
\[
  F^i(D(X^\bullet))\times \ext^j(X^\bullet,Y^\bullet) \to F^{i+j}(D(Y^\bullet)) \text{.}
\]
These pairings, applied to cohomological derived functors, are none
other, by definition, than the Yoneda products. The properties
of these products follow immediately from this definition. We shall not
develop them here.

Let $\cA$, $\cA'$, $\cA''$ be three abelian categories and
$f:\cA\to \cA'$, $g:\cA'\to \cA''$ two additive functors. Suppose that the
functors $\eR^+ f$ and $\eR^+ g$ exist. I do not know how to deduce that the
functor $\eR^+ g\circ f$ exists and even if the latter functor exists, it
is probably not true in general that one has the formula:
\[\xymatrix{
  \eR^+ g\circ f\ar[r]^-\sim
    & \eR^+ g \circ \eR^+ f
}\]
We have however the following proposition: 





\begin{proposition}\label{VIII:4-3-1}
Suppose that the category $\cA'$ has enough
right $g$-acyclic objects and that the category $\cA$ has enough
right $f$-acyclic objects sent by $f$ to right
$g$-acyclic objects. Then the functor $\eR^+ g\circ f$ exists and we have an
isomorphism: 
\[\xymatrix{
  \eR^+ g\circ f \ar[r]^-\sim 
    & \eR^+ g \circ \eR^+ f \text{.}
}\]
\end{proposition}

Analogous statement for left derived functors. 















\section{Examples}\label{VIII:6} % originally chapter 2, section 3










\subsection{The functor \texorpdfstring{$\mathsf{Rhom}^\bullet$}{Hom*}}\label{VIII:6-1}

Let $\cA$ be an abelian category, with enough injectives.

The functor: $Y\mapsto \hom^\bullet(X^\bullet,Y^\bullet)$ (\ref{VIII:1-3})
where $X^\bullet$ is an object of $\eK(\cA)$ and $Y^\bullet$ an object of
$\eK^+(\cA)$, admits a total right derived functor
(Corollary \ref{VIII:4-2-3-1}) which will be denoted $\rhom^\bullet(X^\bullet,-)$.
Let $Y$ be an object of $\eD^+(\cA)$ the functor 
\begin{align*}
  \eK(\cA) &\to \eD(\mathsf{Ab}) \\
  X^\bullet &\mapsto \rhom^\bullet(X^\bullet,Y) 
\end{align*}
is exact and vanishes on acyclic complexes, whence an exact
bifunctor
\[\xymatrix{
  \eD(\cA)^\circ \times \eD^+(\cA) \ar[r]
    & \eD(\mathsf{Ab})
}\]
which we denote by $\mathsf{Rhom}^\bullet$.

When $\cA$ is of finite cohomological dimension, the functor
$\mathsf{Rhom}^\bullet$ extends to $\eD(\cA)^\circ\times \eD(\cA)$
(loc.\ cit.).

If moreover $\cA$ has enough projectives one can derive the
functor $\hom$ in the first argument. The two definitions
coincide in their common domain of validity. 










\subsection{\texorpdfstring{$\Tor$}{Tor} of sheaves}\label{VIII:6-2}

Let $E$ be a ringed site (in particular a ringed topological space),
$\sA$ the sheaf of rings and $\mathsf{Mod}(E)$ the category of
$\sA$-modules. An object $X$ of $\mathsf{Mod}(E)$ is said to be
\emph{$\sA$-flat} if for every exact sequence: 
\[\xymatrix{
  0 \ar[r] 
    & Y' \ar[r] 
    & Y \ar[r] 
    & Y'' \ar[r] 
    & 0 
}\]
the sequence: $0 \to Y'\otimes_\sA X \to Y\otimes_\sA X \to Y''\otimes_\sA X\to 0$
is exact. There are enough $\sA$-flat objects: direct sums
of objects $\sA_U$, where $U$ is an object of $E$ and $\sA_U$ the sheaf
$\sA$ ``extended by zero outside $U$.'' Theorem
\ref{VIII:4-2-3} therefore applies and one can define the functor 
\begin{align*}
  \eD^-(E,\sA) \times \eD^-(E,\sA) &\to \eD^-(E,\sA) \\
  (X,Y) &\mapsto X\lotimes Y 
\end{align*}
total tensor product, left derived functor of the tensor product functor.
One thus defines passing to cohomology the $\Tor_i^\sA(X,Y)$:
local hyper-tor.

The category $\mathsf{Mod}(E)$ having enough injectives, one can
derive the functor: $\hom^\bullet(X,Y)$: complex of local homomorphisms,
whence $\mathsf{Rhom}(X,Y)$ ($X\in \eD$, $Y\in \eD^+$).

Let $X$ and $Y$ be two objects of $\eD^-(E,\sA)$, $Z$ an object of
$\eD^+(E,\sA)$. There exists a functorial isomorphism: 
\[\xymatrix{
  \hom_{\eD(E,\sA)}(X\lotimes Y,Z) \ar[r]^-\sim 
    &\hom_{\eD(E,\sA)}(X,\mathsf{Rhom}(Y,Z)) \text{.}
}\]










\subsection{Left derived functor of inverse image}\label{VIII:6-3}

Let $f:E\to E'$ be a morphism of ringed topological spaces. The
inverse image functor:
\[
  f^\ast:\mathsf{Mod}(E') \to \mathsf{Mod}(E)
\]
is not necessarily exact. But one can define
$f^\ast$-flat objects and there are enough such objects. One can thus
define the functor $\eL^- f^\ast$. One likewise has a duality formula:

Let $X$ be an object of $\eD^-(E,\sA)$ and $Y$ an object of
$\eD^+(E',\sA')$; there exists a functorial isomorphism: 
\[\xymatrix{
  \hom_{\eD(E',\sA')}(\eL^- f^\ast(X),Y) \ar[r]^-\sim 
    & \hom_{\eD(E,\sA)}(X,\eR^+ f_\ast(Y)) \text{.}
}\]
















